\documentclass[border=10pt,crop=true]{standalone}
\usepackage{circuitikz}
% Shared style for 电子学一 Chapter 2 op-amp circuit figures (CircuiTikZ).
\usetikzlibrary{calc}

\ctikzset{
  bipoles/length=0.95cm,
  bipoles/thickness=1,
  resistors/scale=1,
  capacitors/scale=1,
  label distance=2.5mm,
}

\tikzset{
  opfig/.style={font=\small},
  opfigThin/.style={font=\scriptsize},
}

% Geometry (cm) — keep identical across all op-amp schematics
\def\OpInLen{1.55}    % label → input terminal
\def\OpInGap{0.08}    % terminal → wire to op amp +
\def\OpAmpWire{1.00}  % terminal side → op amp + anchor
\def\OpOutLen{1.05}   % op amp out → output terminal
\def\OpJuncL{0.50}    % v_- → feedback junction (left)
\def\OpFbDown{0.55}   % v_- straight down before R_1
\def\OpRvert{1.20}    % R_1 vertical span to ground
\def\OpInSep{0.55}    % dual-input spacing (v_1 / v_2)
\def\OpFbStub{0.40}   % follower: stub right of output before dropping
\def\OpFbClear{0.22}  % follower: below op-amp south for horizontal feedback


\begin{document}
\begin{circuitikz}[american, scale=1.0, transform shape]
  \draw (-\OpInLen, 0) node[left, opfig]{$v_i$} to[short, o-] ++(\OpInGap, 0)
    -- ++(\OpAmpWire, 0)
    node[op amp, noinv input up, anchor=+](OA){};

  % v_- 左侧结点 VM（反馈与反相端相连）
  \draw (OA.-) -- ++(-\OpJuncL, 0) coordinate(VM);

  % 输出端
  \draw (OA.out) to[short, -o] ++(\OpOutLen, 0) node[right, opfig]{$v_o$};

  % 反馈：先右 → 下至三角形下方 → 水平至 VM → 已在 OA.- 处汇合
  \coordinate (FBLOW) at ($(OA.south) + (0, -\OpFbClear)$);
  \draw (OA.out) -- ++(\OpFbStub, 0) coordinate(STUB);
  \draw (STUB) |- (FBLOW) -| (VM);
\end{circuitikz}
\end{document}
